\documentclass[10pt,a4paper]{amsart}
\usepackage[margin=19mm]{geometry}
\usepackage{amsmath,amssymb,mathtools}
\usepackage[scr=rsfs,cal=euler]{mathalfa}
\usepackage{xfrac}
\usepackage[unicode,hidelinks,bookmarks=false]{hyperref}
\newcommand{\ie}{i.e.\ }
\newcommand{\cf}{cf.\ }
\newcommand{\resp}{respectively\ }
\newcommand{\Subsection}{\subsection}
\providecommand{\nobreakdash}{\nobreak\hbox{-}}
\setlength{\emergencystretch}{2em}
\multlinegap0pt
\usepackage{amssymb}
\usepackage{tikz}
\usepackage{enumerate}
\usepackage{bm}
\usepackage{xcolor}
\usepackage[normalem]{ulem}
\newtheorem{proposition}{Proposition}
\newtheorem{lemma}{Lemma}
\newcommand{\MW}{\mathrm{MW}}
\title{Bounds on the Mordell-Weil rank of elliptic fibrations}
\author{Antonella Grassi, Rick Miranda, Kapil Paranjape, Vasudevan Srinivas, Timo Weigand}
\date{Source: arXiv:2603.24666v2 (CC BY 4.0)}
\begin{document}
\maketitle

\noindent\emph{Source and licence.} Bounds on the Mordell-Weil rank of elliptic fibrations. Version arXiv:2603.24666v2 (CC BY 4.0), distributed by arXiv under the Creative Commons Attribution 4.0 International licence, http://creativecommons.org/licenses/by/4.0/, as recorded in the arXiv licence metadata for this exact version and shown on the abstract page https://arxiv.org/abs/2603.24666v2. Full licence notice: \texttt{licenses/arxiv-2603.24666-CC-BY-4.0.txt}.\par\smallskip

\noindent\textit{Editorial scope.} Counted unit: the article's lemma lemma:bsaZ (source0.tex lines 955--962). The article's complete original proof of it is delivered with it (source0.tex lines 965--979, one \texttt{proof} environment of 15 lines). No mathematics is written, repaired or re-derived: the statement and the proof are the article's own lines, in the article's own notation, and no retained line is substituted or re-flowed.  Retained uncounted as the source-local context the proof needs: lines 616--618.  Nothing was omitted from any retained span, so the package carries no omission record.  No retained line contains a citation, so the wrapper carries no bibliography and no bibliography entry.  No retained line contains a figure environment, an image-inclusion command or drawing code, so no image is delivered and the package declares no asset dependency.  Editorial formatting only, in addition: an AMS \texttt{amsart} preamble that restates the article's own declarations and macros used by the retained lines, namely the environments \texttt{proposition}, \texttt{lemma} on separate counters, together with the standard packages those lines need.  The retained environments print in this document as Proposition 1, Lemma 1 in order of appearance, whereas the article prints the same items with the numbering of its own full document.  The complete native source of the article as distributed in the arXiv e-print for this exact version is bundled unchanged as \texttt{sources/197163/source0.tex}.
\begin{proposition}\label{prop:CoxZucker}
  Let $S_a$ be any (rational) section distinct from the zero section $S_0$.  Then there  exist an open set $\mathcal{V} \subset B$  and  an integer $m\in\mathbb{N}$ such that the multiple $m S_a$ (viewed in the Mordell-Weil group) intersects  the singular fibres in the same component as $S_0$ over $\mathcal{V}$.
\end{proposition}

\begin{lemma}\label{lemma:bsaZ}
Let $S_a \in \MW(X/B)$, $ C \subset B $  be a smooth curve contained in $B^{sm}$, the smooth locus of $B$, but such that $ C \not \subset\Sigma_{X/B}$,  and $ Z \stackrel{def}= \pi^{-1}(C)$  . %smooth}. 
Then
\[
 b_{S_a} \cdot C= -{\pi_\ast}_{|\widehat Z}\bigl(\sigma ({S_a}_{|\widehat Z}) \cdot  \sigma({S_a}_{|\widehat Z}) \bigr) = {\langle \sigma ({S_a}_{|\widehat Z}) , \sigma({S_a}_{|\widehat Z})\rangle}_{|_{\widehat Z}}\,
\]
is  well-defined and gives the height pairing $\langle,\rangle_{|_{\widehat Z}}$ for $S_a|_{\widehat Z}\in \MW({\widehat Z}/C)$, where  $ \tau: \widehat Z \to Z$ is a minimal resolution  of $Z$ relative to the morphism $Z \to C$.
\end{lemma}

\begin{proof} 
Recall that $ b_{S_a}
   = -\pi_*   \bigl( \sigma(S_a) \cdot  \sigma(S_a)  \bigr) = \langle  \sigma(S_a),  \sigma(S_a) \rangle.$
  By Proposition \ref{prop:CoxZucker} there exists a positive integer $m$ such that 
$mS_a$ intersects the singular fibers, possibly after restricting to an open neighborhood,  in the same component as $S_0$ and then $ b_{S_a}
    = \frac{1}{m^2}\Big(2 \Lambda + 2\pi_*\bigl((mS_a)\cdot S_0\bigr)\Big)
  $. There exists a Weierstrass model $W_X$ of $X_{|\pi^{-1}(B^{sm})} \to B^{sm}$ and   a Weierstrass model $W_Z$ of $Z\to C$. Now the birational map between $X$ and $W_X$ is an isomorphism in an open set which includes $(mS_a)\cdot S_0$ for a suitable choice of $m$. Likewise, the birational maps between $Z$, $\widehat Z$ and $W_Z$ are isomorphisms in an open set of $((mS_a)\cdot S_0))_{|\pi^{-1}(C)}$. Then 
$\pi_*\bigl((mS_a)\cdot S_0\bigr)\bigr)\cdot C=(\pi_{|\widehat Z})_*\bigl((mS_a)_{|\widehat Z}\cdot {S_0}_{|\widehat Z}\bigr)\Big),$  where  $(mS_a)_{|\widehat Z} \ $ is also the addition in $\widehat Z$, {by base change}. We also have  $K_{\widehat Z}= \pi_{|\widehat Z}^*(K_C+ \Lambda \cdot C)$ and thus
\begin{eqnarray} b_{S_a} \cdot C
   & =& \frac{1}{m^2}\Big(2 \Lambda + 2\pi_*\bigl((mS_a)\cdot S_0\bigr)\Big)\cdot C \nonumber \\
   &=& \frac{1}{m^2}\Big(2 \Lambda \cdot C + 2(\pi_{|\widehat Z})_*\bigl((mS_a)_{|\widehat Z}\cdot {S_0}_{|\widehat Z}\bigr)\Big) \nonumber
\end{eqnarray}
and 
 \[ b_{S_a} \cdot C = -{\pi_\ast}_{|\widehat Z}\bigl(\sigma ({S_a}_{|\widehat Z}) \cdot  \sigma({S_a}_{|\widehat Z}) \bigr) = \langle \sigma ({S_a}_{|\widehat Z}) , \sigma({S_a}_{|\widehat Z})\rangle_{|_{\widehat Z}} \,.
 \]\end{proof}

\end{document}
